
\section{Motivation}
\subsection{What is \emph{random}?}
\begin{itemize}
  \item Some behavior that we cannot perfectly predict, as oppose to
    being \emph{deterministic}.
\end{itemize}

\subsection{Why random?}
\begin{itemize}
\item \textbf{Theologian:}  We can't know God's mind! 
\item \textbf{Copenhagen quantum physicist:} The laws of quantum physics are fundamentally random.
\end{itemize}
  
\subsection{What to do with random?}
Although we cannot perfectly determine a random behavior, we still want to be
able to obtain to some degree its specifications:
\begin{itemize}
  \item \textbf{Probability Theory} is the theory that systematically studies and characterizes randomness.
  \item We will go through an engineering (more operational, less
    about math vigor) treatment of Probability Theory in this course.
\end{itemize}


\subsection{How to study probability?}
\begin{enumerate}
\item \textbf{Intuition}
\begin{itemize}
\item \htmlId{}{<u>Example:</u>} Dr. Wong is a nice guy, he will probably give me a good grade in EEE 5544 .
\item Not systematic, guessing, cannot quantify.
\end{itemize}

\item \textbf{Relative frequency (a.k.a. counting)} 
\begin{itemize}
\item \textbf{\smash{Example:}} Dr. Wong taught EEE 5544 three times in the
  past. There were altogether $200$ students who took Dr. Wong's
  classes. Out of these $200$ students, $50$ of them got the \(A\) grade.
  Thus the chance of getting an \(A\) in Dr. Wong's class is
  \(\frac{50}{200}=\frac{1}{4}\).
\item Assume every grade is ``equally likely'', which may not be the case.
\item What about if there are infinite possibilities?
\end{itemize}

\item \textbf{Axiomatic theory} 
\begin{itemize}
  \item Identify a few basic axioms that fit intuition.
  \item Build a math consistent system based on the axioms.
  \item Utilize measure theory to construct the math system.
  \item Developed by Kolmogorov in the early 1900s.
\end{itemize}
\end{enumerate}
